\scanpage{7}
% archive foliation: 6
\uline{Conséquence} : Si $\widehat{\mathcal{J}}$ est l'idéal d'augm. de $\widehat S$, qu'on filtre $\widehat{\mathcal{J}}$-adiquement, on aura
\[
\operatorname{gr}^1_{\widehat{\mathcal{J}}}(\widehat S)\simeq\Omega
\]
d'où
\[
\operatorname{gr}^{\bullet}_{\widehat{\mathcal{J}}}(\widehat S)\simeq\operatorname{Sym}^{\bullet}(\Omega)
\]
D'autre part, $\mathcal{J}^n/\mathcal{J}^{n+2}$ est une ext. de $\operatorname{gr}^n\simeq\operatorname{Sym}^n(\Omega)$ par $\operatorname{gr}^{n+1}=\operatorname{Sym}^{n+1}(\Omega)$ qu'on va calculer, en notant que l'op. $u_{a*}$ de $a$ sur \deleted{\unsure{$\operatorname{gr}$}} $\widehat{\mathcal{J}}^{(0)n}/\widehat{\mathcal{J}}^{(0)n+2}$ est donnée par
\[
\deleted{\unsure{\sum\omega_1\cdots\omega_n+\sum\omega'_1\cdots\omega'_{n+1}\longrightarrow}}
\]
\[
\begin{aligned}
\underset{\deg n}{f}+\underset{\deg n+1}{g}
&\longrightarrow\frac{f}{(\unsure{1+a|1-a})^n}+\frac{g}{(\unsure{1-a|1+a})^{n+1}}\\
&\unsure{=}f\unsure{-|+}naf+g+\text{termes de degré }\geq n+2
\end{aligned}
\]
Soit donc
\diagram{\begin{tikzcd}[column sep=large]
\Omega\arrow[r,"v_n"]&\underline{\operatorname{Hom}}(\operatorname{Sym}^n\Omega,\operatorname{Sym}^{n+1}\Omega)\arrow[d,phantom,"\simeq"]\\
&\text{Aut d'ext de }\widehat{\mathcal{J}}^{0n}/\widehat{\mathcal{J}}^{0n+2}
\end{tikzcd}}
défini par
\[
v(a)(f)=af.
\]
\deleted{on a alors}
(\uline{NB}. C'est un iso si $\Omega$ inversible). On a alors
\[
\mathcal{T}\times_{\Omega}[\underline{\operatorname{Hom}}(\operatorname{Sym}^n\Omega,\operatorname{Sym}^{n+1}\Omega),-nv_n]
\]
qui donne l'extension.
